\documentclass[12pt,a4paper]{report}

% ── Packages ──────────────────────────────────────────────
\usepackage[margin=1in]{geometry}
\usepackage{graphicx}
\usepackage{booktabs}
\usepackage{array}
\usepackage{longtable}
\usepackage{hyperref}
\usepackage{xcolor}
\usepackage{titlesec}
\usepackage{fancyhdr}
\usepackage{enumitem}
\usepackage{parskip}
\usepackage{setspace}
\usepackage{caption}
\usepackage{listings}
\usepackage{tabularx}
\usepackage{multirow}
\usepackage{amsmath}

% ── Hyperlink colours ─────────────────────────────────────
\hypersetup{
    colorlinks=true,
    linkcolor=blue!70!black,
    urlcolor=blue!70!black,
    citecolor=blue!70!black
}

% ── Code listing style ────────────────────────────────────
\lstset{
    basicstyle=\ttfamily\small,
    backgroundcolor=\color{gray!10},
    frame=single,
    rulecolor=\color{gray!40},
    breaklines=true,
    keywordstyle=\color{blue},
    commentstyle=\color{green!50!black},
    stringstyle=\color{red!70!black},
    numbers=left,
    numberstyle=\tiny\color{gray},
    numbersep=5pt
}

% ── Header / Footer ───────────────────────────────────────
\pagestyle{fancy}
\fancyhf{}
\rhead{\small CSC336 -- Web Technologies}
\lhead{\small Lab Assignment 2}
\rfoot{\thepage}
\lfoot{\small COMSATS University Islamabad, Vehari Campus}
\renewcommand{\headrulewidth}{0.4pt}
\renewcommand{\footrulewidth}{0.4pt}

% ── Section title formatting ──────────────────────────────
\titleformat{\chapter}[display]
  {\normalfont\Large\bfseries\centering}
  {Chapter \thechapter}{12pt}{\large}
\titlespacing*{\chapter}{0pt}{20pt}{15pt}

\titleformat{\section}{\normalfont\large\bfseries}{\thesection}{1em}{}
\titleformat{\subsection}{\normalfont\normalsize\bfseries}{\thesubsection}{1em}{}

% ── Line spacing ──────────────────────────────────────────
\onehalfspacing

% ══════════════════════════════════════════════════════════
\begin{document}

% ── Title Page ────────────────────────────────────────────
\begin{titlepage}
\centering
\vspace*{0.5cm}

{\Large\bfseries COMSATS University Islamabad, Vehari Campus\\[4pt]
Pakistan}\\[0.4cm]

{\normalsize Department of Computer Science}\\[0.6cm]

\rule{\linewidth}{0.5pt}\\[0.3cm]
{\large\bfseries Course: CSC336 -- Web Technologies}\\[0.2cm]
{\large\bfseries Lab Assignment 2}\\[0.3cm]
\rule{\linewidth}{0.5pt}\\[0.6cm]

\textbf{GitHub Repository Link:}\\[0.2cm]
\href{https://github.com/ahsanalich358/Lab-Assignment-2Web.git}{%
\texttt{https://github.com/ahsanali/Web-Tech-Lab-Ass-02}}\\[0.4cm]


\includegraphics[width=0.28\textwidth]{logo.png}\\[1.2cm]

{\large\bfseries
Building a Laravel-Based User Management System\\[4pt]
with Full CRUD Functionality, Profile Image Handling,\\[4pt]
and Email Search Feature}\\[0.8cm]

{\normalsize by}\\[0.3cm]

{\large\bfseries
Ahsan Ali}\\[1.5cm]

\begin{tabular}{ll}
\textbf{Submitted By:} & Ahsan Ali\\[4pt]
\textbf{Student ID:}   & CIIT/SP24-BSE-004/VHR\\[4pt]
\textbf{Instructor:}   & Mam Yasmeen Jana\\[4pt]
\textbf{Semester:}     & Spring 2026\\
\end{tabular}

\vfill
{\small\textit{Department of Computer Science, COMSATS University Islamabad, Vehari Campus}}
\end{titlepage}

% ── Front Matter ──────────────────────────────────────────
\pagenumbering{roman}

% ── Abstract ──────────────────────────────────────────────
\chapter*{Abstract}
\addcontentsline{toc}{chapter}{Abstract}

This report presents the design and implementation of a web-based User Management
System developed using the Laravel PHP framework. The system is built following the
Model--View--Controller (MVC) architectural pattern and provides a complete suite of
CRUD (Create, Read, Update, Delete) operations for managing user records. Key
features include profile image upload and storage, email-based filtering of records, and
input validation with confirmation dialogs to prevent accidental data loss. The
application is deployed locally using Laravel's built-in Artisan development server and
uses MySQL as the backend database. The project demonstrates practical application
of Laravel routing, Eloquent ORM, Blade templating, and form handling in a real-world
web development scenario.

\vspace{1cm}
\textbf{Keywords:} Laravel, CRUD, MVC, PHP, User Management, Image Upload,
Search, Blade Templates, Eloquent ORM, MySQL.

% ── Abbreviations ─────────────────────────────────────────
\chapter*{List of Abbreviations}
\addcontentsline{toc}{chapter}{List of Abbreviations}

\begin{tabularx}{\textwidth}{@{}lX@{}}
\toprule
\textbf{Abbreviation} & \textbf{Full Form} \\
\midrule
CRUD   & Create, Read, Update, Delete \\
MVC    & Model--View--Controller \\
PHP    & Hypertext Preprocessor \\
HTML   & HyperText Markup Language \\
CSS    & Cascading Style Sheets \\
SQL    & Structured Query Language \\
URL    & Uniform Resource Locator \\
ORM    & Object Relational Mapping \\
IDE    & Integrated Development Environment \\
UI     & User Interface \\
DB     & Database \\
HTTP   & Hypertext Transfer Protocol \\
REST   & Representational State Transfer \\
CMD    & Command Prompt / Terminal \\
\bottomrule
\end{tabularx}
% ── Main Content ──────────────────────────────────────────
\clearpage
\pagenumbering{arabic}

% ══════════════════════════════════════════════════════════
\chapter{Introduction and System Overview}
% ══════════════════════════════════════════════════════════

\section{Background and Motivation}

Managing user data is a fundamental requirement in virtually every modern web
application — ranging from academic portals and HR systems to e-commerce platforms
and administrative dashboards. Traditional methods of record keeping, such as
spreadsheets or paper-based forms, are inefficient, error-prone, and difficult to scale.
Web-based systems built on robust frameworks provide a structured, maintainable, and
scalable alternative.

This project, \textbf{Building a Laravel-Based User Management System with Full
CRUD Functionality, Profile Image Handling, and Email Search Feature}, is designed
to address these challenges by implementing a complete user record management
solution using the Laravel framework. Laravel's expressive syntax, built-in ORM
(Eloquent), and Blade templating engine make it an ideal choice for rapid development
of data-driven web applications.

The system enables authorised users to register new records, browse all saved entries
in a structured table, update existing information, and permanently remove records
after explicit confirmation. Profile images can be attached to each user record, and a
real-time email search allows quick retrieval of specific entries from potentially large
datasets.

\section{Tools and Technologies}

Table~\ref{tab:tools} lists the tools and technologies selected for the development of
this project. Each component was chosen based on compatibility, community support,
and suitability for Laravel-based web development.

\begin{table}[h!]
\centering
\caption{Tools and Technologies Used}
\label{tab:tools}
\renewcommand{\arraystretch}{1.3}
\begin{tabularx}{\textwidth}{@{}llX@{}}
\toprule
\textbf{Tool / Technology} & \textbf{Version} & \textbf{Purpose} \\
\midrule
Laravel       & Latest (v11.x)   & Primary PHP framework; implements MVC architecture and handles routing, ORM, and views. \\
PHP           & 8.x              & Server-side scripting language required to execute Laravel applications. \\
Composer      & Latest           & PHP dependency manager used to install and manage Laravel packages. \\
MySQL         & Latest           & Relational database used to persist user records, including image file paths. \\
Blade         & Built-in         & Laravel's templating engine used to build dynamic HTML views and forms. \\
Bootstrap 5   & 5.x              & Frontend CSS framework used to create a responsive, clean user interface. \\
VS Code       & Latest           & Primary code editor used for writing and debugging the application. \\
Git \& GitHub & Latest           & Version control and remote repository hosting for the project codebase. \\
\bottomrule
\end{tabularx}
\end{table}

\section{Project Objectives}

The core objectives guiding the design and development of this system are:

\begin{itemize}[leftmargin=1.5cm]
    \item \textbf{PO-1:} Design and implement a fully functional web application that
    supports Create, Read, Update, and Delete operations on user records using the
    Laravel MVC framework.

    \item \textbf{PO-2:} Extend the base CRUD system with additional usability features
    — specifically, profile image upload and display, email-based record search, and
    input validation with user-confirmation prompts for destructive actions.

    \item \textbf{PO-3:} Follow clean code and architectural conventions (MVC pattern,
    RESTful routing, Eloquent ORM) to produce a maintainable and extensible codebase.
\end{itemize}

% ══════════════════════════════════════════════════════════
\chapter{Implementation}
% ══════════════════════════════════════════════════════════

\section{Development Environment}

The project was built and tested on a local Windows machine using the tools listed in
Table~\ref{tab:devenv}.

\begin{table}[h!]
\centering
\caption{Development Environment Setup}
\label{tab:devenv}
\renewcommand{\arraystretch}{1.3}
\begin{tabularx}{\textwidth}{@{}llX@{}}
\toprule
\textbf{Component} & \textbf{Tool / Technology} & \textbf{Role} \\
\midrule
Backend Framework  & Laravel (v11.x) & MVC architecture, routing, ORM, validation. \\
Programming Lang.  & PHP 8.x         & Server-side execution of Laravel code. \\
Dependency Mgr.    & Composer        & Installs Laravel and all required packages. \\
Database           & MySQL / SQLite  & Stores user records and image file paths. \\
View Engine        & Blade + Bootstrap 5 & Dynamic UI pages: forms, tables, alerts. \\
Code Editor        & Visual Studio Code & Writing, debugging, and managing files. \\
Local Server       & Laravel Artisan & Runs the application via \texttt{php artisan serve}. \\
Version Control    & Git + GitHub    & Tracks code changes and hosts the repository. \\
\bottomrule
\end{tabularx}
\end{table}

\section{Project Setup and Execution}

To run the project on a local machine, execute the following steps in the terminal
from inside the project root directory:

\begin{lstlisting}[language=bash, caption={Project Setup Commands}]
# Step 1 – Install all PHP dependencies
composer install

# Step 2 – Copy environment file and set the database credentials
cp .env.example .env
# Edit .env:  DB_DATABASE=your_db  DB_USERNAME=root  DB_PASSWORD=

# Step 3 – Generate application key
php artisan key:generate

# Step 4 – Run database migrations to create the table
php artisan migrate

# Step 5 – Create the uploads directory (if not present)
mkdir public/uploads

# Step 6 – Start the local development server
php artisan serve
\end{lstlisting}

After the server starts, open \texttt{http://127.0.0.1:8000} in any browser to
access the application.

\section{Core Implementation Details}

\subsection{CRUD Operations}

All CRUD operations are implemented through Laravel's resourceful routing and a
dedicated \texttt{UserController}. The key controller methods are:

\begin{itemize}[leftmargin=1.5cm]
    \item \textbf{index():} Retrieves all user records (with optional email filter)
    and passes them to the \texttt{users/index} Blade view.
    \item \textbf{create():} Returns the \texttt{users/create} registration form view.
    \item \textbf{store():} Validates incoming form data, handles the image upload,
    and persists a new record to the database.
    \item \textbf{edit(\$id):} Retrieves the specified user record and populates the
    edit form view.
    \item \textbf{update(\$id):} Validates and saves changes to the selected user record; replaces the image if a new file is uploaded.

  \includegraphics[width=0.9\textwidth]{1.png}
  
  \includegraphics[width=0.9\textwidth]{2.png}
  
  \includegraphics[width=0.9\textwidth]{3.png}
  
  \includegraphics[width=0.9\textwidth]{4.png}
  
  \includegraphics[width=0.9\textwidth]{5.png}
% ══════════════════════════════════════════════════════════
\chapter*{Conclusion}
\addcontentsline{toc}{chapter}{Conclusion}
% ══════════════════════════════════════════════════════════

This lab assignment successfully demonstrates the practical implementation of a
web-based User Management System using the Laravel PHP framework. The application
covers all four CRUD operations in a clean, maintainable codebase structured
according to the MVC design pattern. Additional features — email-based search,
profile image upload and display, server-side input validation, and JavaScript
confirmation dialogs — extend the system beyond a basic data-entry tool into a
usable, real-world application component.

Key learning outcomes achieved through this project include:
\begin{itemize}[leftmargin=1.5cm]
    \item Hands-on experience with Laravel routing, controllers, models, and
    Blade views.
    \item Practical understanding of Eloquent ORM for database interactions.
    \item Implementation of file upload handling in a PHP/Laravel context.
    \item Application of server-side validation and flash messaging for user feedback.
    \item Use of Git and GitHub for version control throughout the development cycle.
\end{itemize}

The project can be extended in future work by adding user authentication,
role-based access control, pagination, export to CSV/PDF, and deployment to a
cloud hosting platform.
\end{document}
